% =====================================================================
%  preamble.tex — πακέτα, γραμματοσειρές, μορφοποίηση
% =====================================================================

% ---------------- Γλώσσα και γραμματοσειρές (XeLaTeX) ----------------
\usepackage{fontspec}
\usepackage{polyglossia}
\setdefaultlanguage{greek}
\setotherlanguage{english}

% Επιλογή της πρώτης διαθέσιμης γραμματοσειράς με ελληνικούς χαρακτήρες.
\IfFontExistsTF{CMU Serif}{%
  \setmainfont{CMU Serif}%
  \setsansfont{CMU Sans Serif}%
}{%
  \IfFontExistsTF{Linux Libertine O}{%
    \setmainfont{Linux Libertine O}%
    \setsansfont{Linux Biolinum O}%
  }{%
    \setmainfont{DejaVu Serif}%
    \setsansfont{DejaVu Sans}%
  }%
}
\IfFontExistsTF{DejaVu Sans Mono}{%
  \setmonofont[Scale=0.85]{DejaVu Sans Mono}%
}{%
  \setmonofont{CMU Typewriter Text}%
}

% ---------------------------- Σελιδοποίηση ----------------------------
\usepackage[a4paper,top=2.5cm,bottom=2.5cm,inner=3cm,outer=2.5cm,
            headheight=15pt]{geometry}
\usepackage{setspace}
\onehalfspacing
\setlength{\parskip}{0.35em}
\setlength{\parindent}{1.5em}
\emergencystretch=3em

% ----------------------------- Γενικά πακέτα --------------------------
\usepackage{amsmath,amssymb}
\usepackage{graphicx}
\usepackage{xcolor}
\usepackage{booktabs}
\usepackage{longtable}
\usepackage{tabularx}
\usepackage{array}
\usepackage{multirow}
\usepackage{enumitem}
\setlist{itemsep=0.15em,topsep=0.4em}
\usepackage{caption}
\captionsetup{font=small,labelfont=bf,format=hang}
\usepackage{subcaption}
\usepackage{float}
\usepackage{fancyhdr}
\usepackage{titlesec}
\usepackage{tcolorbox}
\tcbuselibrary{skins,breakable}
\usepackage{csquotes}
\usepackage{pdfpages}

% ------------------------------- Γραφικά -------------------------------
\usepackage{tikz}
\usetikzlibrary{arrows.meta,positioning,shapes.geometric,shapes.misc,
                fit,backgrounds,calc,shadows}
\usepackage{pgfplots}
\pgfplotsset{compat=1.17}

% ------------------------------- Χρώματα -------------------------------
\definecolor{critical}{HTML}{B3261E}
\definecolor{warning}{HTML}{C77700}
\definecolor{notice}{HTML}{1E5AA8}
\definecolor{okgreen}{HTML}{2E7D32}
\definecolor{codebg}{HTML}{F6F8FA}
\definecolor{codeframe}{HTML}{D0D7DE}
\definecolor{keyword}{HTML}{8250DF}
\definecolor{comment}{HTML}{6E7781}
\definecolor{string}{HTML}{0A3069}
\definecolor{univblue}{HTML}{1F3A68}

% -------------------------------- Κώδικας -------------------------------
\usepackage{listings}
\lstdefinestyle{thesis}{
  backgroundcolor=\color{codebg},
  frame=single, rulecolor=\color{codeframe},
  basicstyle=\ttfamily\footnotesize,
  keywordstyle=\color{keyword}\bfseries,
  commentstyle=\color{comment}\itshape,
  stringstyle=\color{string},
  numbers=left, numberstyle=\tiny\color{comment}, numbersep=6pt,
  breaklines=true, breakatwhitespace=false,
  showstringspaces=false, tabsize=4,
  columns=fullflexible, keepspaces=true,
  xleftmargin=1.6em, framexleftmargin=1.4em,
  captionpos=b, aboveskip=0.8em, belowskip=0.8em,
  literate={-}{{\char45}}1,
}
\lstset{style=thesis}
\lstdefinelanguage{shell}{
  morekeywords={sudo,python3,pacman,git,cd,echo,tee,touch,modprobe,udevadm,
                systemctl,latexmk,xelatex,biber},
  sensitive=true, morecomment=[l]{\#}, morestring=[b]",
}
\lstdefinelanguage{json}{
  morestring=[b]", stringstyle=\color{string},
  literate={:}{{\textcolor{keyword}{:}}}1 {-}{{\char45}}1,
}
\renewcommand{\lstlistingname}{Απόσπασμα}
\renewcommand{\lstlistlistingname}{Κατάλογος Αποσπασμάτων Κώδικα}

% ------------------------------ Βιβλιογραφία ----------------------------
\usepackage[backend=biber,style=numeric-comp,sorting=none,
            maxbibnames=6,giveninits=true,url=true,doi=true]{biblatex}
\addbibresource{references.bib}
% Χωρίς ελαστικά κενά μέσα στα URL της βιβλιογραφίας.
\biburlbigskip=0mu \biburlnumskip=0mu \biburlucskip=0mu \biburllcskip=0mu
\AtBeginBibliography{\sloppy}

% -------------------------------- Σύνδεσμοι ------------------------------
\usepackage[hidelinks,unicode,bookmarksnumbered]{hyperref}
\hypersetup{
  pdftitle={Probolos: Πύλη Αποδοχής Συσκευών USB με Προεπιλεγμένη Άρνηση},
  pdfsubject={Μεταπτυχιακή Διπλωματική Εργασία},
  pdfkeywords={USB, BadUSB, Linux, ασφάλεια, sysfs, διαχωρισμός προνομίων},
}
\usepackage{bookmark}
\pdfstringdefDisableCommands{%
  \def\en#1{#1}\def\code#1{#1}\def\Probolos{Probolos}%
  \def\sevC{CRITICAL}\def\sevW{WARNING}\def\sevN{NOTICE}}

% ------------------------------ Επικεφαλίδες -----------------------------
\pagestyle{fancy}
\fancyhf{}
\fancyhead[L]{\small\nouppercase{\leftmark}}
\fancyhead[R]{\small\thepage}
\renewcommand{\headrulewidth}{0.4pt}
\fancypagestyle{plain}{\fancyhf{}\fancyfoot[C]{\small\thepage}%
  \renewcommand{\headrulewidth}{0pt}}

\titleformat{\chapter}[display]
  {\normalfont\huge\bfseries\color{univblue}}
  {\large\chaptertitlename\ \thechapter}{12pt}{\Huge}
\titlespacing*{\chapter}{0pt}{-10pt}{30pt}
\titleformat{\section}{\normalfont\Large\bfseries\color{univblue}}{\thesection}{1em}{}
\titleformat{\subsection}{\normalfont\large\bfseries}{\thesubsection}{1em}{}

\setcounter{secnumdepth}{3}
\setcounter{tocdepth}{2}

% ---------------------------- Βοηθητικές εντολές -------------------------
% Κώδικας εντός κειμένου· επιτρέπεται αλλαγή γραμμής μετά από παύλα.
\newcommand{\code}[1]{\texttt{\hyphenchar\font=`\-\relax #1}}
% Διαδρομές αρχείων· αλλαγή γραμμής μετά από «/».
\newcommand{\pth}[1]{\nolinkurl{#1}}
\newcommand{\en}[1]{\textenglish{#1}}
\newcommand{\Probolos}{\textsf{Probolos}}
\newcommand{\sevC}{\textcolor{critical}{\textbf{CRITICAL}}}
\newcommand{\sevW}{\textcolor{warning}{\textbf{WARNING}}}
\newcommand{\sevN}{\textcolor{notice}{\textbf{NOTICE}}}

% Πλαίσιο «σημείωσης» / «ορισμού»
\newtcolorbox{notebox}[1][Σημείωση]{%
  enhanced, breakable, colback=notice!5, colframe=notice!70,
  fonttitle=\bfseries, title=#1, sharp corners=south, boxrule=0.6pt}
\newtcolorbox{warnbox}[1][Προσοχή]{%
  enhanced, breakable, colback=critical!4, colframe=critical!70,
  fonttitle=\bfseries, title=#1, sharp corners=south, boxrule=0.6pt}
\newtcolorbox{defbox}[1][Ορισμός]{%
  enhanced, breakable, colback=okgreen!4, colframe=okgreen!70,
  fonttitle=\bfseries, title=#1, sharp corners=south, boxrule=0.6pt}

% Θεωρήματα / ορισμοί (ελληνική αρίθμηση)
\usepackage{amsthm}
\theoremstyle{definition}
\newtheorem{definition}{Ορισμός}[chapter]
\newtheorem{requirement}{Απαίτηση}[chapter]
\theoremstyle{remark}
\newtheorem*{remark}{Παρατήρηση}

% Αρίθμηση σχημάτων/πινάκων ανά κεφάλαιο
\numberwithin{figure}{chapter}
\numberwithin{table}{chapter}
